\section{Introduction}

Vehicle-to-Everything (V2X) communication enables safety-critical and compute-intensive applications such as cooperative perception, autonomous driving, and real-time hazard avoidance. These applications impose heterogeneous Quality-of-Service (QoS) requirements: safety messages may become unusable within approximately 100 ms, whereas compute and infotainment workloads can tolerate substantially larger delays. At the same time, processing the growing volume of camera, radar, and LiDAR data entirely on resource-constrained vehicle Electronic Control Units (ECUs) is increasingly difficult. Vehicular Edge Computing (VEC) addresses this limitation by offloading computation to roadside Multi-access Edge Computing (MEC) servers.

The effectiveness of offloading, however, depends jointly on where computation is performed and how the offloaded workload reaches the edge. No single Radio Access Technology (RAT) is uniformly optimal across task types and mobility conditions. Long-Term Evolution (LTE), ITS-G5 (the European Intelligent Transport Systems band at 5.9 GHz), WiFi-6, and 5G New Radio (5G-NR) exhibit different latency, throughput, coverage, and mobility characteristics. A RAT suitable for a delay-tolerant compute task may therefore be unsuitable for a safety-critical task generated by the same vehicle. Moreover, binary local/offload decisions cannot exploit parallel local and MEC processing. These observations motivate treating RAT selection and partial computation offloading as a joint decision rather than as independent optimization stages.

A second challenge emerges as vehicle density increases. Vehicles associated with the same Roadside Unit (RSU) compete not only for wireless resources but also for a shared MEC processing queue. Consequently, aggressive offloading that performs well for an isolated vehicle may degrade when many vehicles simultaneously submit workloads. This degradation can also be hidden by conventional communication metrics. Packet Reception Ratio (PRR) counts a packet as successfully delivered even when it arrives after its application deadline; for safety-critical traffic, however, a late result may be operationally equivalent to a lost one. Evaluation must therefore distinguish successful reception from deadline-compliant delivery.

To address these coupled challenges, we propose MA-ROPE, a multi-agent framework that jointly learns RAT selection and partial computation offloading under shared MEC contention. The core insight is that RAT choice and offload ratio are mutually dependent: the selected RAT determines transmission latency and reliability, which in turn determines the feasible offload ratio; conversely, the offload ratio sets transmission duration, which affects queue contention and RSU delay for all co-located vehicles. Optimizing these decisions sequentially yields suboptimal policies because downstream interactions are ignored.

Each vehicle selects among LTE, ITS-G5, WiFi-6, and 5G-NR while adapting its offloaded workload fraction $\rho \in \{0.25, 0.50, 0.75, 1.0\}$ to channel conditions, task requirements, energy budget, and current RSU queue state. The MDP state explicitly includes shared queue depth $Q$, enabling each vehicle's policy to account for the decisions of co-located vehicles. A Dueling Double Deep Q-Network (D3QN) decomposes value and advantage streams, enabling stable learning under multi-vehicle interactions. Federated Averaging (FedAvg) coordinates policies across vehicle clusters without requiring vehicles to share raw observations, preserving operational privacy while enabling cooperative learning at the policy level.

Accordingly, this work investigates whether a federated multi-RAT offloading policy can preserve deadline-compliant V2X delivery as fleet density and shared MEC contention increase.

Despite extensive research on VEC offloading, existing approaches generally address only subsets of heterogeneous RAT selection, partial computation offloading, multi-vehicle MEC contention, mobility-aware channel behavior, and federated cooperative learning. Recent distributed deep RL methods (e.g., \cite{zhao2024}, \cite{guo2026}) employ federated learning or game-theoretic coordination, but typically optimize RAT selection or offload ratio independently, rather than jointly. Bandit-based approaches (e.g., \cite{wu2021}, \cite{haq2026}) adapt network selection to density, but do not model partial execution or shared queue dynamics. Task-splitting strategies (e.g., \cite{jing2019}) optimize fairness or latency, but rely on analytical optimization rather than learning from channel variability and queue state. No prior work jointly optimizes RAT selection, partial-offload ratio, and shared RSU queue contention within a federated deep RL framework that explicitly models deadline compliance across multiple fleet densities.

MA-ROPE addresses this gap by bringing these dimensions together and introducing a new evaluation discipline: not only performance at a fixed operating point, but critically, how well deadline-compliant service survives increasing fleet contention. We quantify this via the Deadline Retention Ratio (DRR), which measures how much of a policy's reference-density on-time performance remains at higher densities. This metric reveals that aggressive offloading policies that perform well at low vehicle density can collapse catastrophically as the fleet grows, a failure mode invisible to scalar operating-point metrics.

The main contributions of this work are:


First, we formulate RAT selection and partial computation offloading as a unified MDP decision, not as sequential stages. Unlike prior work that selects the RAT first and then optimizes offload ratio (or vice versa), we model the action space as (RAT, $\rho$) pairs, directly encoding the mutual dependence. Critically, the MDP state includes the shared RSU queue depth, enabling each vehicle to account for others' decisions and avoid the "AllOffload collapse" observed when vehicles naively offload all traffic without considering queue saturation.
Second, We coordinate D3QN training across vehicle clusters via Federated Averaging (FedAvg), enabling cooperative learning without centrally pooling raw vehicle observations. The key innovation is that the per-agent state explicitly includes the shared RSU queue depth, making each vehicle's learned policy aware of co-located vehicle demand. This differs from pure distributed RL by incorporating shared resource constraints at the agent level, not requiring a central coordinator during execution.
Third, We distinguish communication reliability (PRR) from application-useful on-time delivery (OTR). We introduce the DRR, which normalizes a policy's on-time delivery at a target density relative to its own reference-density performance, enabling fair comparison across heterogeneous policies without dependence on utility-function constants.
Finally, We evaluate MA-ROPE and four baselines across multiple fleet sizes ($N \in \{8, 16, 32\}$) under URLLC traffic. We demonstrate that while baseline policies suffer catastrophic failure modes at high contention, the proposed federated MDP formulation robustly maintains deadline compliance, link reliability, and low computational overhead across varying densities.

